\documentclass[a4paper]{article}
\usepackage{amsfonts}
\usepackage{amsmath}
\usepackage{amssymb}
\usepackage{graphicx}     

\begin{document}
  
\noindent \large Auteur: Sabawoon Enayat \\
\noindent \large Studierichting: Informatica\\
\noindent \large Oefeningenreeks 4A nr. 42.2 \\

\newcommand{\dint}{\displaystyle\int}

\medskip

\normalsize

Opgave: $I = \dint \dfrac{dx}{x^2 - 4}$\\

$= \dint \dfrac{dx}{(x - 2)(x + 2)}$\\

$\Rightarrow \dfrac{1}{(x - 2)(x + 2)} = \dfrac{A}{x - 2} + \dfrac{B}{x + 2}$\\

$A = \left. \dfrac{1}{x + 2} \right|_{x=2} = \dfrac{1}{4}$\\

$B = \left. \dfrac{1}{x - 2} \right|_{x=-2} = -\dfrac{1}{4}$\\

$I = \dfrac{1}{4}\left(\dint \dfrac{dx}{x - 2} - \dint \dfrac{dx}{x + 2}\right)$\\

$= \dfrac{1}{4}\left(\dint \dfrac{d(x - 2)}{x - 2} - \dint \dfrac{d(x + 2)}{x + 2}\right)$\\

$= \dfrac{1}{4}\left(\ln |x - 2| - \ln |x + 2|\right) + c$\\

$= \dfrac{1}{4} \ln \left|\dfrac{x - 2}{x + 2}\right| + c$\\

\begin{thebibliography}{999}
%\bibitem{a} Referentie
\end{thebibliography}
\end{document} 
